\section{Spectral Targets in Affine-Invariant Geometry}
\label{sec:information-rdgc}

We begin with affine-invariant SPD geometry and isolate the static spectral
principles used by the structured families.

\subsection{Relative SPD geometry}

Let
\[
  f(\theta)=\frac12\theta^\top H\theta,
  \qquad H\in\SPD^d.
\]
For a metric \(G\in\SPD^d\), gradient flow is
\(\dot\theta=-G^{-1}H\theta\).  Define the relative metric
\[
  S=H^{-1/2}GH^{-1/2}.
\]
Then \(\kappa(G^{-1}H)=\kappa(S)\), where the left side denotes the ratio of
generalized eigenvalues in \(Hv=\lambda Gv\).

The affine-invariant metric has line element and distance
\[
  \norm{Z}_{S,\AI}
  =\fro{S^{-1/2}ZS^{-1/2}},
  \qquad
  d_{\AI}(S_0,S_1)
  =\fro{\log(S_0^{-1/2}S_1S_0^{-1/2})}.
\]
The SPD cone is a finite-dimensional Hadamard manifold
\cite{bhatia2007positive,pennec2006riemannian}.

For \(K\ge1\), set
\[
  \Ck=\{S\in\SPD^d:\kappa(S)\le K\}.
\]

\begin{theorem}[Full-SPD spectral benchmark]
\label{thm:full-spd-benchmark}
Let \(y_1\le\cdots\le y_d\) be the ordered log-eigenvalues of
\(S_0\in\SPD^d\).  Then
\[
  D_K(S_0)
  :=d_{\AI}(S_0,\Ck)
  =
  \min_{c\in\R}
  \left(
  \sum_{i=1}^d\dist(y_i,[c,c+\log K])^2
  \right)^{1/2}.
\]
\end{theorem}

Thus the unconstrained benchmark clips the relative log-spectrum into an
interval of width \(\log K\).  Restricted families constrain which spectral
and eigenspace motions are admissible.

\subsection{Structured metric families and target distance}

Let \(\F\subset\SPD^d\) be an admissible metric family and define its
Hessian-relative realization
\[
  \mathcal S_\F(H)
  =
  \{H^{-1/2}GH^{-1/2}:G\in\F\}.
\]
A realization specifies a path-connected component, treatment of scale
freedom, an admissible absolutely continuous path class, and the induced or
quotient affine-invariant length.  Smooth embedded, closed geodesically convex,
quotient, and stratified families therefore remain distinct cases.

\begin{definition}[Structured affine-invariant target distance]
\label{def:restricted-complexity}
For \(G_0\in\F\) and \(S_0=H^{-1/2}G_0H^{-1/2}\), define
\[
\begin{aligned}
  D_{K,\F}(S_0;H)
  =
  \inf\{&\Len_{\AI}(S_{[0,T]}):
  S(0)=S_0,\ S(T)\in\Ck,\\
  &S_t\in\mathcal S_\F(H),
  \ S_{[0,T]}\text{ admissible}\}.
\end{aligned}
\]
The value is \(+\infty\) if the relative target is empty or lies in a
different admissible component.
\end{definition}

This target distance records the least intrinsic affine-invariant motion from the initial geometry to a structured condition target.  The endpoint records visible
preconditioning quality, and the admissible curve is eliminated by the
infimum.  For closed geodesically convex realizations it reduces to a unique
target projection; nonconvex or stratified families retain the variational
definition without inheriting that uniqueness.

\subsection{Generating geodesically convex information}

For \(H_a\in\SPD^d\) and a convex permutation-invariant function
\(\phi_a:\R^d\to\R\), define
\[
  \Omega_a(G)
  =
  \phi_a\!\left(
  \log\lambda(G^{-1/2}H_aG^{-1/2})
  \right).
\]

\begin{theorem}[Information-spectrum generation law]
\label{thm:structured-log-spectral-mother}
Let \(\F\subset\SPD^d\) be nonempty, closed, and geodesically convex.  If
\(\rho:\R^q\to\R\) is convex and nondecreasing in every coordinate, then
\[
  \mathcal J(G)
  =
  \rho(\Omega_1(G),\ldots,\Omega_q(G))
\]
is closed and geodesically convex on \(\F\).  The conclusion holds pointwise
for continuously time-varying \(H_{t,a},\phi_{t,a},\rho_t\).

If every \(\phi_a(x+c\mathbf1)=\phi_a(x)\), then \(\mathcal J(cG)=\mathcal
J(G)\).  If \(\phi_a\) is \(L_a\)-Lipschitz and \(\rho\) is
\(L_\rho\)-Lipschitz, replacing \(H_a\) by \(\widehat H_a\) changes the
objective by at most
\[
  \sup_{G\in\F}
  |\widehat{\mathcal J}(G)-\mathcal J(G)|
  \le
  L_\rho
  \left(\sum_{a=1}^qL_a^2
  d_{\AI}(H_a,\widehat H_a)^2\right)^{1/2}.
\]
\end{theorem}

The condition-number objective is obtained from
\[
  \phi_{\mathrm{osc}}(x)=\max_i x_i-\min_i x_i,
  \qquad
  \mathcal J_H(G)=\log\kappa(G^{-1}H).
\]

\subsection{Static action value, exact penalty, and contraction}

The static theorem is stated on an abstract finite-dimensional Hadamard
manifold so that it can be reused by every closed geodesically convex family.

\begin{theorem}[Restricted target projection and exact penalty]
\label{thm:restricted-geometrodynamic-mother}
Let \((\mathcal M,d)\) be a finite-dimensional Hadamard manifold,
\(\F\subset\mathcal M\) a nonempty closed geodesically convex set, and
\(\mathcal J:\F\to\R\cup\{+\infty\}\) proper, closed, geodesically convex,
and bounded below.  Suppose
\[
  \mathcal C_\tau=\{G\in\F:\mathcal J(G)\le\tau\}
\]
is nonempty.  For quadratic kinetic action and terminal indicator
\(\iota_{\mathcal C_\tau}\), the value is
\[
  V(t,G)=\frac{d(G,\mathcal C_\tau)^2}{2(T-t)}.
\]
For \(G_0\in\F\), the least admissible length is
\[
  \mathfrak D_{\tau,\F}^{\mathcal J}(G_0)
  =d(G_0,\mathcal C_\tau)
  =d(G_0,P_\tau),
  \qquad
  P_\tau=P_{\mathcal C_\tau}(G_0).
\]
The projection and constant-speed minimizing geodesic are unique, and for
every \(Y\in\mathcal C_\tau\),
\[
  d(G_0,Y)^2
  \ge
  d(G_0,P_\tau)^2+d(P_\tau,Y)^2.
\]

If families or thresholds are nested, their distances are monotone.  If the
objectives satisfy
\[
  \sup_{G\in\F}|\widehat{\mathcal J}(G)-\mathcal J(G)|\le\varepsilon,
\]
then
\[
  \mathfrak D_{\tau+\varepsilon,\F}^{\mathcal J}(G_0)
  \le
  \mathfrak D_{\tau,\F}^{\widehat{\mathcal J}}(G_0)
  \le
  \mathfrak D_{\tau-\varepsilon,\F}^{\mathcal J}(G_0)
\]
whenever the outer values are finite.

Assume a strict feasible geometry \(G_s\) satisfies
\(\mathcal J(G_s)\le\tau-\sigma\), \(\sigma>0\), and put
\(R=d(G_0,G_s)>0\).  On
\(\mathcal D=\F\cap\overline B_R(G_0)\), every
\[
  \Lambda>\frac{2R^2}{\sigma}
\]
makes
\[
  \Phi_\Lambda(G)
  =
  \frac12d(G_0,G)^2
  +\Lambda\pospart{\mathcal J(G)-\tau}
\]
an exact penalty with unique minimizer \(P_\tau\).  For every \(\eta>0\),
the proximal iteration
\[
  G_{r+1}
  =
  \argminop_{G\in\mathcal D}
  \left\{\Phi_\Lambda(G)+\frac{d(G,G_r)^2}{2\eta}\right\}
\]
satisfies
\[
  d(G_r,P_\tau)^2
  \le
  (1+2\eta)^{-r}d(G_0,P_\tau)^2.
\]
If \(R=0\), then \(G_0=P_\tau\).
\end{theorem}

The proximal recursion is a variational convergence benchmark: its rate is
conditional on solving each global proximal subproblem.  No claim is made that
those subproblems are cheaper than the original target projection for an
arbitrary family.

For the condition-number target, take \(\mathcal J=\mathcal J_H\) and \(\tau=\log K\).  Then
\(\mathfrak D_{\tau,\F}^{\mathcal J}=D_{K,\F}\) for every closed
geodesically convex realization.

\subsection{Normal response and the nearest-versus-best question}

Let \(\F\) now be a nonempty, closed, geodesically convex smooth embedded
submanifold.  In whitened coordinates define
\[
  \widehat T_G\F
  =\{G^{-1/2}UG^{-1/2}:U\in T_G\F\},
  \qquad
  \widehat N_G\F=(\widehat T_G\F)^\perp.
\]
For \(S_a=G^{-1/2}H_aG^{-1/2}\), let
\[
  \mathfrak Z_a(G)
  =
  \partial_X[\phi_a(\lambda(X))]_{X=\log S_a}.
\]
Here \(\partial_X\) is the Euclidean convex subdifferential in the whitened
symmetric-matrix coordinate \(X\); the affine-invariant tangent/cotangent
identification is applied only after this whitening.

\begin{lemma}[Log-spectral pullback identity]
\label{lem:log-spectral-pullback}
Let \(S\in\SPD^d\), let \(\phi:\R^d\to\R\) be convex and
permutation invariant, and set \(F(Y)=\phi(\lambda(Y))\) on
\(\mathbb S^d\).  For every \(X\in\mathbb S^d\),
\[
  \frac{\dd}{\dd t}\bigg|_{t=0+}
  F\!\left(\log(e^{-tX/2}Se^{-tX/2})\right)
  =
  \max_{Z\in\partial F(\log S)}-\tr(ZX).
\]
Consequently, the whitened AIRM subdifferential of the corresponding relative
log-spectral objective is exactly \(-\partial F(\log S)\), including at
repeated eigenvalues.
\end{lemma}

\begin{theorem}[Normal-response law]
\label{thm:normal-response-law}
A point \(G\in\F\) globally minimizes the information objective in
\cref{thm:structured-log-spectral-mother} if and only if there exist
\[
  \theta\in\partial\rho(\Omega_1(G),\ldots,\Omega_q(G)),
  \qquad
  Z_a\in\mathfrak Z_a(G),
\]
such that
\[
  \sum_{a=1}^q\theta_aZ_a\in\widehat N_G\F.
\]
Thus optimality is exact cancellation of the tangent information force.

For a single \(H\), let \(P=P_\F(H)\) be its unique nearest affine-invariant point
in \(\F\) and set \(S=P^{-1/2}HP^{-1/2}\).  Then
\(\log S\in\widehat N_P\F\), and \(P\) also minimizes the log-spectral
difficulty \(\Omega_{\phi,H}\) if and only if
\[
  \mathfrak Z_\phi(P)\cap\widehat N_P\F\ne\emptyset.
\]
\end{theorem}

\begin{corollary}[Sharp scale-closed threshold]
\label{thm:scale-closed-threshold}
Let \(d\ge2\) and let \(\F\subset\SPD^d\) be nonempty and closed under
positive rescaling.  Define
\[
  K_\F^\star(H)=\inf_{G\in\F}\kappa(G^{-1}H),
  \qquad
  \delta_\F(H)=d_{\AI}(H,\F),
  \qquad
  \alpha_d=\sqrt{\frac{d}{\lfloor d^2/4\rfloor}}.
\]
Then
\[
  \log K_\F^\star(H)\ge\alpha_d\delta_\F(H),
  \qquad
  K_\F^\star(H)\ge\exp\{\alpha_d d_{\AI}(H,\F)\}.
\]
The constant \(\alpha_d\) is sharp over general scale-closed families.
\end{corollary}

\begin{proposition}[Monotonicity and full-SPD lower bound]
\label{prop:monotonicity}
If \(\F_1\subset\F_2\) and every admissible \(\F_1\)-path is an admissible
\(\F_2\)-path with the same length, then
\[
  D_{K,\F_1}(S_0;H)\ge D_{K,\F_2}(S_0;H).
\]
In particular, \(D_{K,\F}(S_0;H)\ge D_K(S_0)\) whenever restricted paths are
ambient SPD paths.  Also
\(D_{K_1,\F}\ge D_{K_2,\F}\) for \(K_1\le K_2\).
\end{proposition}

\begin{proposition}[Intrinsic submanifold distance]
\label{prop:submanifold-distance}
If \(\mathcal S_\F(H)\) is a smooth embedded submanifold equipped with the
ambient induced length, then
\[
  D_{K,\F}(S_0;H)
  =
  d_{\mathcal S_\F(H)}
  (S_0,\mathcal S_\F(H)\cap\Ck).
\]
The distance is \(+\infty\) when the target lies outside the admissible
path-connected component.
\end{proposition}
